\element{Conservation}{
  \begin{question}{conservationCharge}\bareme{1}
    L'équation locale de conservation de la charge s'écrit
    \begin{multicols}{2}
      \begin{reponses}
        \bonne{$\frac{\partial \rho}{\partial t}=-div\vec{j}$}
        \mauvaise{$\frac{\partial \rho}{\partial t}=\vec{grad}\vec{j}$}
        \mauvaise{$\frac{\partial \vec{j}}{\partial x}=\vec{grad}\vec{\rho}$}
        \mauvaise{$\frac{\partial \vec{j}}{\partial x}=div\vec{\rho}$}
      \end{reponses}
    \end{multicols}
  \end{question}
}

\element{unites}{
  \begin{questionmult}{uniteJ}\bareme{haut=2,p=0}
    Le vecteur densité de courant électrique se mesure en
    \begin{multicols}{3}
      \begin{reponses}
        \bonne{\SI{}{A.m^{-2}}}
        \mauvaise{\SI{}{A.m^{-3}}}
        \mauvaise{\SI{}{A.m^2}}
        \mauvaise{\SI{}{A.m^3}}
        \mauvaise{\SI{}{C.s^{-1}}}
        \mauvaise{\SI{}{C.m^{-1}}}
      \end{reponses}
    \end{multicols}
  \end{questionmult}
}
\element{unites}{
  \begin{questionmult}{uniteRho}\bareme{haut=2,p=0}
    La densité volumique de charge se mesure en
    \begin{multicols}{4}
      \begin{reponses}
        \bonne{\SI{}{C.m^{-3}}}
        \mauvaise{\SI{}{A.m^3}}
        \mauvaise{\SI{}{A.s^{-1}.m^{-3}}}
        \mauvaise{\SI{}{C.m^{3}}}
      \end{reponses}
    \end{multicols}
  \end{questionmult}
}
\setgroupmode{unites}{cyclic}

\element{relations}{
  \begin{questionmult}{relations}\bareme{haut=2,p=0}
    Parmi ces relations, lesquelles sont correctes lorsque tous les porteurs de charge sont identiques ?
    \begin{multicols}{4}
      \begin{reponses}
        \bonne{$\vv{j}=nq\vv{v}$}
        \bonne{$\vv{j}=\rho_\text{libre}\vv{v}$}
        \bonne{$\rho=nq$}
        \mauvaise{$\vv{j}=n\vv{v}$}
        \mauvaise{$n=\rho q$}
        \mauvaise{$\rho_\text{libre}=-\rho$}
      \end{reponses}
    \end{multicols}
  \end{questionmult}
}
\setgroupmode{relations}{cyclic}

\element{Ohm}{
  \begin{question}{ohmLocale}\bareme{1}
    La loi d'Ohm locale s'écrit s'écrit (avec $\gamma$ la conductivité électrique)
    \begin{multicols}{2}
      \begin{reponses}
        \mauvaise{$\vec{E}=\gamma\vec{j}$ avec $\gamma$ en \SI{}{\ohm.m^{-1}}}
        \bonne{$\vec{j}=\gamma\vec{E}$ avec $\gamma$ en \SI{}{S.m^{-1}}}
        \mauvaise{$\vec{j}=\gamma\vec{E}$ avec $\gamma$ en \SI{}{S.m}}
        \mauvaise{$\vec{E}=\gamma\vec{j}$ avec $\gamma$ en \SI{}{\ohm.m^}}
      \end{reponses}
    \end{multicols}
  \end{question}
}

\element{puissance}{
  \begin{questionmult}{puissanceJouleLocale}\bareme{haut=2,p=0}
    La puissance volumique dissipée par effet Joule
    \begin{multicols}{2}
      \begin{reponses}
        \bonne{Se transforme en chaleur}
        \bonne{S'écrit $\vec{j}.\vec{E}$}
        \bonne{Se mesure en \SI{}{W.m^{-3}}}
        \mauvaise{N'existe que dans les isolants}
        \mauvaise{Est nulle dans les métaux}
      \end{reponses}
    \end{multicols}
  \end{questionmult}
}

\element{densite}{
  \begin{question}{densiteCu}\bareme{1}
    Le cuivre a pour masse volumique $\mu=\SI{9}{g.cm^{-3}}$ et comme masse molaire $M=\SI{60}{g.mol^{-1}}$. On donne $\mathcal{N}_a=\SI{6e23}{u.s.i.}$. Sans calculatrice, calculer la densité particulaire.
    \begin{multicols}{4}
      \begin{reponses}
        \bonne{\SI{9e28}{u.s.i.}}
        \mauvaise{\SI{9e25}{u.s.i.}}
        \mauvaise{\SI{3e26}{u.s.i.}}
        \mauvaise{\SI{4e18}{u.s.i.}}
        \mauvaise{\SI{9e-22}{u.s.i.}}
        \mauvaise{\SI{1e-29}{u.s.i.}}
        \mauvaise{\SI{3e27}{u.s.i.}}
        \mauvaise{\SI{3e28}{u.s.i.}}
      \end{reponses}
    \end{multicols}
  \end{question}
}
\element{densite}{
  \begin{question}{densiteSi}\bareme{1}
    Le silicium a pour masse volumique $\mu=\SI{2}{g.cm^{-3}}$ et comme masse molaire $M=\SI{30}{g.mol^{-1}}$. On donne $\mathcal{N}_a=\SI{6e23}{u.s.i.}$. Sans calculatrice, calculer la densité particulaire.
    \begin{multicols}{4}
      \begin{reponses}
        \bonne{\SI{4e28}{u.s.i.}}
        \mauvaise{\SI{4e25}{u.s.i.}}
        \mauvaise{\SI{9e18}{u.s.i.}}
        \mauvaise{\SI{1e22}{u.s.i.}}
        \mauvaise{\SI{1e-22}{u.s.i.}}
        \mauvaise{\SI{3e-29}{u.s.i.}}
        \mauvaise{\SI{1e-19}{u.s.i.}}
        \mauvaise{\SI{3e-26}{u.s.i.}}
      \end{reponses}
    \end{multicols}
  \end{question}
}
\setgroupmode{densite}{cyclic}